\documentclass[11pt]{amsart}

\usepackage[T1]{fontenc}
\usepackage[utf8]{inputenc}
\usepackage{lmodern}
\input{glyphtounicode}
\pdfgentounicode=1

\usepackage{amsmath,amssymb}
\usepackage{array,booktabs,tabularx}
\usepackage{graphicx}
\usepackage{geometry}
\usepackage{xcolor}
\usepackage{hyperref}
\usepackage{enumitem}
\geometry{margin=1in}
\setlength{\emergencystretch}{3em}
\hypersetup{
  colorlinks=true,
  linkcolor=blue!45!black,
  citecolor=blue!45!black,
  urlcolor=blue!45!black,
  pdftitle={A Numerical Study of Cyclic Multilinear Maps, Orthogonal Projectors, and Multiresolution Tensor Representations},
  pdfauthor={Eliuth Chavero Jasso},
  pdfsubject={Fourteen numerical experiments on projectors, symmetries, subspace transport and tensor compression},
  pdfkeywords={numerical linear algebra, principal angles, orthogonal projector, associator, cyclic symmetry, higher-order SVD, gauge invariance, identifiability}
}

\newcommand{\norm}[1]{\lVert#1\rVert}
\newcommand{\ran}{\operatorname{ran}}

\newtheorem{observation}{Observation}[section]
\theoremstyle{remark}
\newtheorem{remark}[observation]{Remark}

\title[A numerical study of cyclic multilinear maps and projectors]{A Numerical Study of Cyclic Multilinear Maps,\\ Orthogonal Projectors, and Multiresolution\\ Tensor Representations}
\author{Eliuth Chavero Jasso}
\address{Independent Researcher, Apizaco, Tlaxcala, Mexico}
\thanks{None of the results below has been reproduced by anyone other than the author, and
no assessment of originality has been carried out. The finite-dimensional theory of
projected composition trees \cite{ChaveroJassoTrees} is a separate study and is not used
here; see Remark~\ref{rem:separatestudy}.}
\date{31 July 2026}

\begin{document}
\maketitle

\begin{abstract}
We report fourteen numerical experiments on a family of cyclically symmetric multilinear
maps equipped with orthogonal projectors onto learned subspaces, together with a swept
extension of four of them. Four results are exact identities that hold by construction and
are reported as such: idempotence and self-adjointness of $UU^{*}$, cyclic invariance
after exact symmetrisation, agreement of two independent implementations of a reduced
tensor, and invariance of a corrected objective under a change of orthonormal basis. Four
are negative results in the tested regime: a parameterised model of a commutator does not
outperform the zero predictor in any of the fifteen trained checkpoints available; three
independently trained resolutions exhibit no transport of learned subspaces, with
principal angles close to $\pi/2$ throughout; the corresponding tensor factors do not
persist across those resolutions; and three seeds reaching comparable training loss
converge to subspaces that are pairwise nearly orthogonal. Two are methodological
corrections found by the study's own controls: comparing orthonormal frames by
unconstrained unitary Procrustes is vacuous, because the unitary group acts transitively
on orthonormal frames of equal size; and an objective intended to be invariant under a
change of basis was not, until it was summed over all columns. The remaining results are
exploratory. A swept extension over $208$ distinct configurations, each executed on a CPU
and on a GPU, finds the GPU between $3.2$ and $3.5$ times slower than the CPU throughout
$n\in\{12,24,48,96\}$, and finds the observed associator ratio decreasing monotonically
with ambient dimension, from $0.40$ at $n=12$ to $0.12$ at $n=96$, against a triangle bound
of $2$.
\end{abstract}

\subjclass[2020]{65F99, 65Y20, 15A69, 62-08}
\keywords{numerical linear algebra, principal angles, orthogonal projector, associator,
cyclic symmetry, higher-order singular value decomposition, gauge invariance,
identifiability, negative results}

\section{Experimental questions and study design}
\label{sec:design}

\subsection{The object under study}

Fix an ambient dimension $n$, a reduced rank $r$ and an arity $a$. The objects of the
study are: a cyclically symmetric $a$-linear map $\mu$ on $\mathbb C^{n}$, stored in a
rank-$R$ canonical polyadic representation \cite{KoldaBader2009}; an isometry
$U\in\mathbb C^{n\times r}$ with $U^{*}U=I_r$, and the associated orthogonal projector
$P=UU^{*}$; and a reduced tensor obtained by contracting $\mu$ against $U$ in every slot.
Parameters are fitted by gradient descent against objectives built from a closure defect,
an associator defect and a commutator residual, all defined in
Section~\ref{sec:definitions}.

\subsection{The fourteen questions}

Each experiment answers one question. We list them here, and each subsequent section
answers one. The parenthetical labels are the identifiers used in the accompanying
software \cite{ChaveroJassoSoftware} and are given once, here, so that each result can be
traced to its implementation; they are not used again.

\begin{enumerate}[leftmargin=*,label=(Q\arabic*)]
\item Are the projector identities satisfied numerically? (implementation label: block A)
\item Can a parameterised model reproduce a commutator? (block B)
\item How do finite nested-commutator norms behave across families of projectors? (block C)
\item Is thresholded spectral projection stable under perturbation? (block D)
\item Do learned subspaces transport across resolutions? (block E)
\item Is the fitted subspace locally identifiable, and invariant under a change of basis? (block F)
\item How large is the closure defect of a fitted map? (block G)
\item How large is the associator, against its triangle bound? (block H)
\item Do two independent implementations of the reduced tensor agree? (block I)
\item How should tensors at different resolutions be compared? (block J)
\item Is a fitted reduced tensor more compressible than a random one? (block K)
\item How much basis freedom remains after canonicalisation? (block L)
\item Do tensor factors persist across resolutions? (block M)
\item Does cyclic symmetrisation produce a cyclic map, and how large is a generalised
Jacobi residual? (block N)
\end{enumerate}

\subsection{Reporting conventions}

Three conventions are used throughout, and departures from them are the main defects we
found in an earlier version of this work.

\begin{description}[leftmargin=*]
\item[Identities by construction are labelled as such.] Several quantities vanish because
of how they are defined, not because of anything that was fitted. Where this is the case we
say so, and we do not report the vanishing as evidence about the fitted object.
\item[A configuration is not an execution.] A \emph{configuration} is one distinct
mathematical setting: a choice of arity, dimension, rank, representation rank and seed. An
\emph{execution} is one run of one configuration on one device. Executions of the same
configuration on different devices, and executions resumed from a checkpoint, are not
independent observations and are never counted as such.
\item[Statistical language is used only when a sampling model is stated.] Where no
sampling population, statistic, uncertainty calculation and inferential procedure are
specified, we write ``observed over $N$ sampled configurations'' and not ``statistically
validated''.
\end{description}

\begin{remark}[This study is separate from the projected-tree theory]
\label{rem:separatestudy}
The companion article \cite{ChaveroJassoTrees} proves upper bounds for the error incurred
by recursively projecting the intermediate outputs of a multilinear composition tree, and
leaves the exact constants open. That study and this one share vocabulary --- projectors,
associators, multilinear maps --- but not objects: no experiment below involves a
composition tree or a recursively projected evaluation, and no numerical result below bears
on the constants of \cite{ChaveroJassoTrees}. In particular, the triangle bound of value
$2$ appearing in Section~\ref{sec:associator} is a bound on the norm of an associator of a
single map, and is unrelated to the coefficient $k-1$ of \cite{ChaveroJassoTrees}, which
bounds the projected error of a signed combination of trees. The numeral is the same by
coincidence.
\end{remark}

\section{Definitions}
\label{sec:definitions}

Let $\mu$ be an $a$-linear map on $\mathbb C^{n}$, $U\in\mathbb C^{n\times r}$ an isometry
and $P=UU^{*}$.

\emph{Cyclic symmetrisation.} Let $\sigma$ rotate the input slots and
$\Pi_{\mathrm{cyc}}=\frac1a\sum_{j=0}^{a-1}\sigma^{j}$. Since $\sigma$ acts unitarily on
the coefficient tensor, $\Pi_{\mathrm{cyc}}$ is an orthogonal projector, and
$\mu_{\mathrm{cyc}}=\Pi_{\mathrm{cyc}}\mu$ is cyclic. The \emph{cyclic defect} of $\mu$ is
$\norm{\mu-\Pi_{\mathrm{cyc}}\mu}_F$.

\emph{Closure defect.} The normalised sampled closure defect is
\[
 \mathcal E_{\mathrm{closure}}
 =\mathbb E\left[\frac{\norm{(I-P)\mu(PX_1,\ldots,PX_a)}^{2}}
 {\norm{\mu(PX_1,\ldots,PX_a)}^{2}+\varepsilon}\right],
\]
the $X_i$ being independent standard complex Gaussian vectors. This is a sampled average,
not the operator norm of the closure-leakage map, and it does not bound that operator norm.

\emph{Associator.} With $A(x,y,z)=(x\circ y)\circ z-x\circ(y\circ z)$ for the binary
product induced by anchoring one slot, the triangle inequality gives
$\norm{A(x,y,z)}\le2\widehat M^{2}\norm{x}\norm{y}\norm{z}$, where $\widehat M$ is a
numerical estimate of the operator norm of $\circ$. The \emph{associator ratio} is the
observed left-hand side divided by the right-hand side.

\emph{Commutator residual.} For a Hermitian $\Delta$, the identity
$[\Delta,P]=K\Delta P-P\Delta K$ holds with $K=I$, and both sides have rank at most $2r$.
The parameterised model is $C_\theta=U\Phi U^{*}\Delta K-K\Delta U\Phi^{*}U^{*}$ with
$\Phi$ a free $r\times r$ matrix. The \emph{unexplained relative residual} is
$\norm{[\Delta,P]-C_\theta}_F^{2}/\norm{[\Delta,P]}_F^{2}$; the zero predictor
$C_\theta=0$ gives the value $1$.

\emph{Principal angles.} Subspaces are compared by principal angles
\cite{BjorckGolub1973,GolubVanLoan2013}; $\pi/2\approx1.5708$ is complete orthogonality.

\section{Orthogonal projector identities}
\label{sec:projector}

For any isometry $U$, the matrix $P=UU^{*}$ satisfies $P^{2}=P$ and $P^{*}=P$ identically.
Numerically, over five seeds and over all $416$ executions of the swept extension of
Section~\ref{sec:sweep}, the idempotence residual $\norm{P^{2}-P}_F/\norm{P}_F$ stayed
below $3.3\times10^{-16}$ and the self-adjointness residual below $7.7\times10^{-17}$; an
exact $n=2$ case agrees to a tolerance of $10^{-14}$.

\begin{remark}
\label{rem:byconstruction}
These identities hold by construction. They confirm that the implementation forms $UU^{*}$
correctly and that the arithmetic is well conditioned. They are not evidence that the
subspace $\ran P$ is relevant to anything, and we never use them as such.
\end{remark}

\section{Approximation of a commutator by a parameterised model}
\label{sec:commutator}

\subsection{The model is not capacity-limited in isolation}

Training $U$ against the unexplained residual alone drives that residual to approximately
$10^{-9}$ across five seeds, matching a closed-form optimal-$\Phi$ ceiling. The model can
therefore represent the target when nothing else competes.

\subsection{It nevertheless fails in every trained checkpoint}

In every one of the fifteen historical checkpoints carrying the relevant data, the model
performed at or worse than the zero predictor.

\begin{observation}
\label{obs:commutator}
The proposed approximation did not outperform the zero approximation in any of the fifteen
available trained models.
\end{observation}

This conclusion is restricted to the parameterisation, the objectives and the trained
models examined here. It is not an impossibility statement for the model class.

\subsection{Which factor is responsible}

Table~\ref{tab:ablation} reports seven training regimes, one run each, at $n=16$, $r=4$,
representation rank $4$, and $400$ optimisation steps.

\begin{table}[h]
\centering\small
\caption{Seven training regimes, one run per regime. ``Unexplained'' is the relative
residual defined in Section~\ref{sec:definitions}; the zero predictor gives $1$. Values are
reproduced from the stored experiment record.}
\label{tab:ablation}
\begin{tabular}{@{}lrrr@{}}
\toprule
regime & unexplained & closure & associator \\
\midrule
commutator objective alone & $8.10\times10^{-5}$ & $0.310$ & $464.1$ \\
\quad with closure objective & $3.51\times10^{-4}$ & $4.11\times10^{-6}$ & $74.6$ \\
\quad with associator objective & $8.49\times10^{-3}$ & $0.111$ & $4.77\times10^{-6}$ \\
all objectives jointly & $6.00\times10^{-3}$ & $4.79\times10^{-4}$ & $1.07\times10^{-6}$ \\
map frozen, projector trained & $0.9925$ & $1.56\times10^{-2}$ & $0.233$ \\
projector frozen, map trained & $2.46\times10^{-7}$ & $0.738$ & $580.9$ \\
staged: competing objectives, then commutator & $4.65\times10^{-4}$ & $0.217$ & $141.1$ \\
\bottomrule
\end{tabular}
\end{table}

Two rows carry the interpretation. Freezing the map and training only the projector leaves
the unexplained residual at $0.9925$, barely below the zero predictor. Freezing the
projector and training only the map --- with $U$ left at its random initial value ---
reaches $2.46\times10^{-7}$, which is the numerical floor of that run and not exactly zero.
Adding the associator objective to the commutator objective degrades the residual by about
two orders of magnitude relative to training the commutator objective alone, whereas adding
the closure objective does not.

\begin{observation}
\label{obs:mechanism}
In these seven runs, the explanatory capacity resides in the parameters of the map rather
than in the fitted subspace, and the commutator objective conflicts with the associator
objective but not with the closure objective.
\end{observation}

This is a single run per regime. It is an exploratory result, not a controlled experiment
with replicates, and it is offered as the best-supported reading of the seven runs rather
than as an established mechanism.

\begin{remark}[Relation to the underspecification literature]
\label{rem:underspecification}
A pipeline is \emph{underspecified} when it returns many distinct predictors with
equivalently strong test performance which nonetheless behave differently in deployment
\cite{DAmourEtAl2022}. The finding of Section~\ref{sec:identifiability} below --- three
seeds reaching comparable loss but converging to nearly orthogonal subspaces --- is
conceptually related to that phenomenon. The finding of the present section is different in
mechanism: it concerns a conflict between two objectives within a single training run,
rather than a multiplicity of equally good solutions. We record the connection as a
conceptual one; the model, the objective and the conclusion analysed here are not those of
\cite{DAmourEtAl2022}, and this remark is not an assessment of originality.
\end{remark}

\section{Finite nested-commutator diagnostics}
\label{sec:nested}

We report Frobenius norms of finite nested commutators for four families of projectors:
adversarially optimised, random, smooth, and spatially localised. The observed norms were
$4.35$, $3.84$, $2.85$ and $2.00$ respectively.

That the adversarially optimised family produces the \emph{largest} norms is the point of
the experiment: it shows that the reported figures are not a favourable selection.

\begin{remark}[Scope]
\label{rem:beals}
This experiment computes finite nested-commutator norms and nothing else. It is not a test
of membership in a pseudodifferential class and carries no microlocal-regularity content.
Such a test would require symbol estimates of the kind recalled in the companion analytic
article, and none is performed here.
\end{remark}

\section{Stability of thresholded spectral projectors}
\label{sec:snapping}

Given a Hermitian matrix with a spectral gap, thresholding its spectrum yields a projector
onto the corresponding invariant subspace. Perturbing by $\epsilon H$ with $H$ Hermitian,
rank recovery and a Davis--Kahan-type bound \cite{DavisKahan1970} hold up to the point at
which the gap closes. At $\epsilon=1.2$ the gap closes from $1.0$ to $0.51$ and the rank is
misrecovered.

\begin{observation}
\label{obs:snapping}
The tested perturbations confirm that the spectral-gap condition is necessary: rank
recovery fails exactly when the gap closes.
\end{observation}

The experiment used genuine near-projectors, not only exact ones.

\section{Transport of learned subspaces across resolutions}
\label{sec:transport}

Three resolutions, $n=12$, $18$ and $24$, were independently initialised and independently
trained. For each ordered pair, one frozen Gaussian-kernel lift operator maps the subspace
of the lower resolution into the higher one, and the transported subspace is compared with
the independently learned subspace there by principal angles. Two baselines are used: a
random subspace of the same dimension, and a subspace obtained by interpolation.

\begin{table}[h]
\centering\small
\caption{Largest principal angle in radians, by resolution pair. Complete orthogonality is
$\pi/2\approx1.571$. The trained lift beats both baselines in two of the three forward
pairs.}
\label{tab:transport}
\begin{tabular}{@{}lrrrr@{}}
\toprule
pair & trained lift & random baseline & interpolation baseline & backward \\
\midrule
$12\to18$ & $1.410$ & $1.480$ & $1.527$ & $1.331$ \\
$12\to24$ & $1.472$ & $1.416$ & $1.468$ & $1.540$ \\
$18\to24$ & $1.407$ & $1.553$ & $1.446$ & $1.370$ \\
\bottomrule
\end{tabular}
\end{table}

All quantities reported in Table~\ref{tab:transport} lie between $1.331$ and $1.553$
radians, that is, between $76.3^{\circ}$ and $89.0^{\circ}$. The transported angles
themselves lie between $1.407$ and $1.472$. The trained lift improves on both baselines in
two of the three forward pairs, and the margins where it does are $0.02$ to $0.15$ radians,
small relative to how close every condition already is to complete orthogonality.

\begin{observation}
\label{obs:transport}
The tested multiresolution models provided no evidence of transport of the learned subspace
across the examined resolutions.
\end{observation}

The scope is three resolutions, one lift operator per pair, one training run per
resolution. This is a negative result in the tested regime and not an impossibility
statement.

\section{Local identifiability and invariance under a change of basis}
\label{sec:identifiability}

\subsection{An objective that was not invariant, and its correction}

Any objective built from $U$ should be invariant under $U\mapsto UV$ for unitary $V$, since
$UV$ spans the same subspace. The first version of the test objective used a single column
of $U$ and was therefore not invariant: the test detected a discrepancy of $0.87\,\%$ under
such a change of basis. Summing over all $r$ columns produces a Frobenius-norm quantity,
which is exactly invariant under right multiplication by a unitary. This is an identity,
and after the correction the discrepancy is zero to machine precision.

\begin{remark}
\label{rem:owncontrol}
The defect was found by the experiment's own required invariance control, not by inspection.
Reporting it is the reason the control exists.
\end{remark}

\subsection{Curvature along the gauge direction}

The exact Hessian and the generalised Gauss--Newton approximation $2J^{\top}J$ are computed
separately and never conflated. The generalised Gauss--Newton matrix is positive
semidefinite by construction; the exact Hessian is indefinite at a non-minimum, as expected.
Along the infinitesimal direction generated by $U\mapsto UV$, the Hessian curvature is
$-3.5\times10^{-18}$, numerically flat, which is what invariance predicts.

\subsection{Different seeds, nearly orthogonal subspaces}

Three seeds all reach comparable near-zero training loss. The largest pairwise principal
angle between the three converged subspaces is $1.55$ radians, approximately $89^{\circ}$.

\begin{observation}
\label{obs:identifiability}
Over the three seeds tested, comparable training loss was attained at subspaces that are
pairwise nearly orthogonal. The subspace is not identified by the objective in the tested
regime.
\end{observation}

Three seeds are three observations, not a statistical sample, and no interval is reported.

\section{Approximate closure of multilinear maps}
\label{sec:closure}

The closure defect of a fitted map was evaluated over $2000$ sampled configurations of the
input vectors, and separately maximised by gradient ascent. The adversarial search found
defects at least as large as the largest random sample in every case.

\begin{observation}
\label{obs:closure}
The closure defect was observed over $2000$ sampled inputs, with an adversarial maximum at
least as large as the largest sampled value.
\end{observation}

\begin{remark}[No sampling model is claimed]
\label{rem:nosampling}
We do not describe this as a statistically validated result. Doing so would require stating
the sampling population, the statistic, the number of independent observations, the
uncertainty calculation, the inferential procedure and the limits of generalisation, and
this study specifies none of them. No interval-arithmetic or sum-of-squares bound on the
closure defect was derived.
\end{remark}

\section{Bounds for associators}
\label{sec:associator}

The triangle bound $\norm{A(x,y,z)}\le2\widehat M^{2}\norm{x}\norm{y}\norm{z}$ was never
violated over $500$ trials at a single configuration ($300$ random inputs and $200$
adversarially optimised). The largest observed ratio was $0.452$. The mean cosine between
the two terms of the associator was $0.067$: the terms are close to uncorrelated, rather
than anti-aligned as the worst case of the triangle inequality would require.

\begin{observation}
\label{obs:associator}
Across the configurations tested, the observed associator ratio remained well below the
triangle bound, with a maximum of $0.957$ over the swept extension of
Section~\ref{sec:sweep}. No analytic improvement of the constant $2$ is derived here, so
the exact constant is undetermined.
\end{observation}

Section~\ref{sec:sweep} extends this measurement over a grid and finds a pattern not
visible at a single configuration.

\section{Extraction of reduced tensors}
\label{sec:reduced}

Two independent implementations of the reduced tensor --- one by explicit loops, one by a
contraction exploiting the low-rank representation, including the cyclic average --- agree
to better than $10^{-10}$ in double precision and $10^{-4}$ in single precision. An exact
rational case at $n=2$, $r=1$ matches a hand computation exactly, with no floating-point
arithmetic involved.

This experiment tests extraction correctness only. It says nothing about whether the
extracted tensor is meaningful.

\section{Gauge-invariant comparison of tensors}
\label{sec:gauge}

Comparing two orthonormal frames $U_1,U_2\in\mathbb C^{n\times r}$ by solving the
unconstrained unitary Procrustes problem \cite{Schonemann1966} over all of $\mathrm U(n)$ is
vacuous: the unitary group acts transitively on orthonormal $r$-frames in $\mathbb C^{n}$,
so a unitary carrying $U_1$ to $U_2$ always exists, whatever the two frames are.

\begin{observation}
\label{obs:procrustes}
An early version of the comparison tool used exactly this procedure and consequently
reported independently drawn random tensors as equivalent. The defect was found by the
experiment's own required negative control, and was corrected by comparing subspaces with
principal angles \cite{BjorckGolub1973} instead.
\end{observation}

The corrected methodology reports four quantities separately --- raw, Procrustes-aligned
within the admissible group, permutation-aligned, and amplitude ratio --- rather than
collapsing them into one score. This is an improvement over comparing a single eigenbasis
heuristic, and it is the methodology used in Sections~\ref{sec:transport}
and~\ref{sec:persistence}.

\section{Multilinear singular-value compression}
\label{sec:hosvd}

Using the multilinear singular value decomposition \cite{DeLathauwerEtAl2000}, a fitted
reduced tensor requires mode ranks $[3,4,4]$ out of a full rank of $4$ to capture $99\,\%$
of its energy, against $[4,4,4]$ for a random tensor matched in norm.

\begin{observation}
\label{obs:hosvd}
The fitted tensor is more compressible than the norm-matched random null in one of three
modes, and equally compressible in the other two.
\end{observation}

This is a modest difference at a single configuration, and a single random tensor is not a
null distribution.

\section{Canonical representatives and residual basis freedom}
\label{sec:canonical}

After canonicalisation, the residual basis freedom is reported explicitly by the size of
each spectral cluster. A deliberately constructed degenerate case, with eigenvalues
$2,1,1$, is detected. Applying an arbitrary --- not infinitesimal --- unitary rotation
within the two-dimensional degenerate eigenspace produces a frame that differs pointwise
(Frobenius difference above $0.1$) while spanning the same subspace (largest principal angle
below $10^{-8}$).

This is the concrete content of the statement that canonicalisation determines a
representative only modulo the residual freedom.

\section{Persistence of tensor factors across resolutions}
\label{sec:persistence}

Applying the corrected comparison of Section~\ref{sec:gauge} to the three independently
trained resolutions of Section~\ref{sec:transport}: the mode ranks required are inconsistent
across the three resolutions ($\{3,3,4\}$, $\{3,4,4\}$, $\{3,3,3\}$), and the mean pairwise
principal angle between corresponding modes is $1.01$ radians.

\begin{observation}
\label{obs:persistence}
The tested multiresolution models provided no evidence of persistent tensor factors across
the examined resolutions.
\end{observation}

One mode of the nine compared shows a near-exact match, $1.5\times10^{-8}$ radians. We have
no explanation for it. It is recorded here as an unresolved anomaly rather than absorbed
into the summary above.

\section{Cyclic symmetrisation and a generalised Jacobi residual}
\label{sec:cyclic}

The cyclic defect before symmetrisation is $4.60$; after applying $\Pi_{\mathrm{cyc}}$ it is
$8.2\times10^{-33}$, thirty-one orders of magnitude smaller.

\begin{remark}[This is an identity]
\label{rem:cyclicidentity}
$\Pi_{\mathrm{cyc}}\mu$ is cyclic by construction, because $\Pi_{\mathrm{cyc}}$ is the
orthogonal projector onto the cyclically invariant subspace. The post-symmetrisation value
is at the numerical floor because the identity is exact, not because a symmetry was learned.
The figure is reported to confirm the implementation, and for no other purpose.
\end{remark}

One exact, versioned generalised Jacobi expression --- the full antisymmetrisation of the
associator over the six permutations of three arguments --- is implemented twice
independently; the two implementations agree to $2.8\times10^{-16}$, and a deliberate sign
mutation is detected. The supremum of the corresponding ratio is not shown to be bounded:
the adversarial maximum is $5.98$ at a single configuration against a mean of $0.43$, and
$4.65$ to $5.92$ across the $48$ configurations of the pilot sweep. This ratio is
undetermined.

The relation of this expression to Filippov's fundamental identity \cite{Filippov1985} and
to other conventions for generalised Jacobi identities has not been checked.

\section{Swept extension over a grid}
\label{sec:sweep}

The experiments above use a single configuration or a few seeds. Four of them ---
Sections~\ref{sec:projector}, \ref{sec:closure}, \ref{sec:associator}
and~\ref{sec:cyclic} --- were extended over a grid, on the hardware described in
\cite{ChaveroJassoSoftware}.

\subsection{Design and accounting}

\begin{table}[h]
\centering\small
\caption{Design of the swept extension. Each configuration was executed once on the CPU and
once on the GPU. The two executions of a configuration are not independent observations.}
\label{tab:sweepdesign}
\begin{tabular}{@{}lrrl@{}}
\toprule
stage & configurations & executions & grid \\
\midrule
pilot & $48$ & $96$ & $a\in\{3,4\}$, $n\in\{12,24\}$, $r\in\{3,6\}$, $R\in\{4,8\}$, $3$ seeds \\
broad & $160$ & $320$ & as above with $n\in\{12,24,48,96\}$, $5$ seeds \\
\midrule
total & $208$ & $416$ & \\
\bottomrule
\end{tabular}
\end{table}

There were no failures in either stage. Summed over the per-execution records, the wall
time was $942.8$ seconds for the pilot and $6300.8$ seconds for the broad stage; the
scheduler reported $945.5$ and $6316.5$ seconds respectively, the difference being
scheduling overhead outside the timed region.

Two further stages --- adaptive intensification of boundary configurations, and reruns with
held-out seeds in extended precision with deterministic algorithms --- were specified but
not executed. Sections~\ref{sec:commutator}, \ref{sec:nested}, \ref{sec:snapping},
\ref{sec:transport}, \ref{sec:identifiability}, \ref{sec:reduced}, \ref{sec:gauge},
\ref{sec:hosvd}, \ref{sec:canonical} and~\ref{sec:persistence} were not extended over the
grid and remain single-configuration evidence.

\subsection{Device throughput}

\begin{table}[h]
\centering\small
\caption{Mean wall time per execution, in seconds, over the $160$ configurations of the
broad stage. Each mean is over $40$ executions.}
\label{tab:throughput}
\begin{tabular}{@{}rrrr@{}}
\toprule
$n$ & CPU & GPU & ratio \\
\midrule
$12$ & $9.44$ & $30.64$ & $3.24$ \\
$24$ & $9.31$ & $30.80$ & $3.31$ \\
$48$ & $9.83$ & $34.50$ & $3.51$ \\
$96$ & $7.90$ & $25.09$ & $3.17$ \\
\bottomrule
\end{tabular}
\end{table}

\begin{observation}
\label{obs:throughput}
Over the tested range $n\in\{12,24,48,96\}$, the GPU was between $3.17$ and $3.51$ times
slower than the CPU for these four experiments, with no trend towards crossover.
\end{observation}

The ratio is stable across an eightfold range in ambient dimension, so the effect is not a
small-$n$ artefact. The practical consequence for further work on these four experiments,
on this hardware, is to use the CPU. No asymptotic-complexity statement and no comparison
between hardware platforms follows from this measurement.

\subsection{The associator ratio decreases with ambient dimension}

Over the $96$ executions of the pilot the observed associator ratio ranged from $0.130$ to
$0.957$ with mean $0.334$; over the $320$ executions of the broad stage it ranged from
$0.042$ to $0.957$ with mean $0.240$. Broken down by ambient dimension, the mean ratio was
$0.404$ at $n=12$, $0.257$ at $n=24$, $0.182$ at $n=48$ and $0.118$ at $n=96$.

\begin{observation}
\label{obs:associatortrend}
The mean observed associator ratio decreases monotonically with ambient dimension over the
tested range, moving further from the triangle bound as $n$ grows.
\end{observation}

The single-configuration value $0.452$ of Section~\ref{sec:associator}, measured at $n=16$,
is consistent with this trend rather than an outlier.

\subsection{Reproducibility of the generalised Jacobi ratio}

The adversarial maximum of the generalised Jacobi ratio reproduced between $4.65$ and
$5.92$ over the $48$ pilot configurations, against $5.98$ at the original single
configuration. The finding is stable across configurations.

\section{Provenance of the historical evidence}
\label{sec:provenance}

An earlier body of evidence for these questions was re-examined before this study was
carried out, and the re-examination changed what that evidence supports.

\begin{observation}
\label{obs:lineage}
The nineteen distinct historical runs are all exploratory: every one records an evaluation
mode of \texttt{screening}, including the run directories whose names suggest otherwise.
The nine runs recorded in the run log share a single seed and form one continuous chain of
non-strict resumptions rooted at a single bootstrap; they are not nine independent trials.
Two distinct script hashes appear across the chain, so the script changed during the
campaign without this being recorded.
\end{observation}

The scoring procedure used at the time computed a global score as the mean of per-experiment
scores, with an inconclusive experiment contributing half credit and an experiment that
never ran contributing full credit; no failing state existed. A mean over such a vocabulary
cannot decrease when an experiment fails to run, so we do not use it. Throughout this
article each experiment is reported separately and no global score is formed.

\section{Limitations}
\label{sec:limitations}

\begin{enumerate}[leftmargin=*]
\item The negative results of Sections~\ref{sec:commutator}, \ref{sec:transport},
\ref{sec:identifiability} and~\ref{sec:persistence} are established for the objectives
actually tested. For tractability these used a closure-only or two-term objective rather
than the full weighted objective of the historical campaign, and
Section~\ref{sec:commutator} itself shows that the choice of objective changes the outcome
materially.
\item The swept extension covers four of the fourteen experiments. The remaining ten are
single-configuration evidence and were not run on the GPU at all.
\item Two further sweep stages were specified and not executed.
\item Single random tensors are used as nulls in Section~\ref{sec:hosvd}. A single draw is
not a null distribution.
\item Timing measurements come from one machine and one software stack.
\item No result here has been reproduced by anyone other than the author.
\item No assessment of originality has been carried out. A bounded literature search
located directly relevant prior work for three of this study's own methodological choices
--- a protocol of the same conservative design in an unrelated domain, established
provenance tooling for run-lineage reconstruction, and the underspecification literature
\cite{DAmourEtAl2022} for the phenomenon of Section~\ref{sec:identifiability}. None of the
methodological corrections described above should be read as claimed novel.
\end{enumerate}

\section{Summary}
\label{sec:summary}

\begin{table}[h]
\centering\small
\caption{Summary of the fourteen experiments and the swept extension.}
\label{tab:summary}
\begin{tabularx}{\linewidth}{@{}lXl@{}}
\toprule
Question & Conclusion & Evidence \\
\midrule
Projector identities (\S\ref{sec:projector}) & satisfied & identity by construction \\
Commutator model (\S\ref{sec:commutator}) & not supported in the tested models & negative result \\
Nested commutators (\S\ref{sec:nested}) & adversarial family gives the largest norms & exploratory \\
Spectral projector stability (\S\ref{sec:snapping}) & gap condition is necessary & counterexample \\
Subspace transport (\S\ref{sec:transport}) & not observed & negative result \\
Identifiability (\S\ref{sec:identifiability}) & subspace not identified & negative result \\
Basis invariance (\S\ref{sec:identifiability}) & holds after correction & identity by construction \\
Closure defect (\S\ref{sec:closure}) & measured & observed over $2000$ samples \\
Associator bound (\S\ref{sec:associator}) & constant undetermined & exploratory \\
Reduced tensor (\S\ref{sec:reduced}) & implementations agree & exact check \\
Tensor comparison (\S\ref{sec:gauge}) & unconstrained Procrustes is vacuous & counterexample \\
Compressibility (\S\ref{sec:hosvd}) & modest, in one mode & exploratory \\
Residual basis freedom (\S\ref{sec:canonical}) & quantified & identity by construction \\
Factor persistence (\S\ref{sec:persistence}) & not observed; one anomaly & negative result \\
Cyclic symmetrisation (\S\ref{sec:cyclic}) & identity & identity by construction \\
Generalised Jacobi ratio (\S\ref{sec:cyclic}) & undetermined & exploratory \\
Device throughput (\S\ref{sec:sweep}) & CPU faster throughout & negative result \\
Associator trend (\S\ref{sec:sweep}) & decreases with $n$ & exploratory \\
\bottomrule
\end{tabularx}
\end{table}

Four of the questions received negative answers in the tested regime, two produced
methodological corrections that were found by the study's own controls, four are identities
that hold by construction, and the remainder are exploratory. No global score is formed from
these, and no result is described as established beyond the regime in which it was measured.
Submission of this study is deferred pending independent reproduction and an assessment of
originality, neither of which has been carried out.

\bibliographystyle{amsplain}
\bibliography{references}

\end{document}
